% Image credits for the photographs and AI illustrations of Book 4 (University, Year 2), French edition.
% Machine-readable license records live in images/book4/CREDITS.md; keep this
% file in sync with frontmatter/image-credits-book4.tex (one \item per item there).
% Every credit, author, dataset and licence is kept exactly as in the English file.
\chapter*{Crédits des images}

Les dessins, cartes et graphiques de ce livre sont originaux~; les cartes
et graphiques tracés à partir de données publiées en citent la source dans
la légende ou dans cette liste. Les photographies sont utilisées sous leurs
licences libres respectives, avec nos remerciements à leurs auteurs~; les
liens complets vers les sources et les adresses des licences de chaque
photographie figurent dans \texttt{images/book4/CREDITS.md}, dans le dépôt
du livre.

\bigskip
À côté des photographies et des dessins originaux, certaines figures sont
des illustrations générées par IA, produites avec le générateur d'images
d'OpenAI à partir de consignes rédigées par les éditeurs du livre, et
vérifiées pour leur exactitude scientifique avant d'être retenues. Elles
ne montrent que des scènes génériques, jamais un lieu réel, une personne
réelle ou un spécimen réel. La consigne complète de chaque image générée
est consignée dans \texttt{images/book4/ai/PROMPTS.md}, dans le dépôt.

\bigskip
\noindent Photographies et données, par chapitre~:
\begin{itemize}\small
  \item Chap.~1 (le milieu interstellaire et la formation des étoiles)~: le disque protoplanétaire de HL Tauri, ALMA (ESO/NAOJ/NRAO), CC BY 4.0~; les Piliers de la création dans la nébuleuse de l'Aigle, NASA, ESA and the Hubble Heritage Team (STScI/AURA), domaine public.
  \item Chap.~2 (formation des planètes et météorites)~: la chondrite carbonée d'Allende et les figures de Widmanstätten de la météorite de fer de Gibeon, James St.~John, CC BY 2.0.
  \item Chap.~3 (les éléments de la Terre)~: xénolite mantellique de lherzolite dans un basalte, Kilbourne Hole, James St.~John, CC BY 2.0.
  \item Chap.~5 (magmas et mise en place)~: dyke près de Shiprock, Dexter Perkins, CC0 (Wikimedia Commons).
  \item Chap.~6 (éruptions volcaniques)~: fontaine de lave sur le cône du Pu`u~`\=O`\=o, Kīlauea, J.\,D. Griggs, USGS, domaine public~; la colonne éruptive du mont Saint Helens, 18 mai 1980, Austin Post, USGS, domaine public~; la colonne éruptive du Pinatubo, 12 juin 1991, Dave Harlow, USGS, domaine public.
  \item Chap.~7 (fusion du manteau)~: échantillon de dunite, Mike Davis (GeoDIL), CC0.
  \item Chap.~8 (différenciation magmatique)~: niveaux de chromitite du complexe du Bushveld, Kevin Walsh, CC BY 2.0~; gabbros lités de l'intrusion de Skaergaard, C. Tegner, CC BY 4.0.
  \item Chap.~10 (réactions métamorphiques et faciès)~: schiste bleu à lawsonite, Ward Creek, Californie, et éclogite, Almenningen, Norvège, James St.~John, CC BY 2.0.
  \item Chap.~11 (thermobarométrie et chemins P--T--t)~: micaschiste à grenat, James St.~John, CC BY 2.0~; grain de coésite dans une éclogite, en lumière polarisée analysée, J.\,R. Smyth, CC BY-SA 1.0.
  \item Chap.~12 (fusion de la croûte~: migmatites et granites)~: affleurement de migmatite à Nashua, New Hampshire, Kevin M. Gill, CC BY 4.0~; filon de leucogranite, Anne Burgess (Geograph), CC BY-SA 2.0.
  \item Chap.~13 (anatomie d'une chaîne de montagnes~: les Alpes)~: laves en coussins de l'ophiolite du Chenaillet, Lionel Duchoiselle, CC BY-SA 3.0~; le Cervin vu de Zermatt, Ximonic (Simo Räsänen), CC BY-SA 4.0.
  \item Chap.~14 (chaînes de collision et chaînes andines comparées)~: le plateau tibétain entre le Népal et Shigatsé, Jan Kranendonk, CC BY-SA 4.0~; l'Altiplano bolivien, Luca Galuzzi, CC BY-SA 2.5.
  \item Chap.~15 (massifs anciens et cycle de Wilson)~: Siccar Point, falaise sud-est, dave souza, CC BY-SA 4.0.
  \item Chap.~16 (bassins sédimentaires comparés)~: plate-forme de forage semi-submersible en mer du Nord, Erik Christensen, CC BY-SA 3.0.
  \item Chap.~17 (îles océaniques)~: l'atoll de Kure vu de la Station spatiale internationale, NASA, domaine public.
  \item Chap.~18 (la Terre précambrienne)~: komatiite à texture spinifex à Pyke Hill, Ontario, James St.~John, CC BY 2.0~; carbonate de couverture de la formation de Nuccaleena, Australie-Méridionale, James St.~John, CC BY 2.0.
  \item Chap.~19 (la Terre phanérozoïque)~: piliers de craie de Monument Rocks, Kansas, James St.~John, CC BY 2.0~; le niveau Bonarelli, Picenus, CC BY-SA 4.0.
  \item Chap.~20 (fossiles et paléoenvironnements)~: \emph{Waptia} et \emph{Marrella} des schistes de Burgess, James St.~John, CC BY 2.0~; le spécimen de Berlin d'\emph{Archaeopteryx}, H.~Raab, CC BY-SA 3.0~; \emph{Zoophycos} à Sagres, Amada44, CC BY 4.0.
  \item Chap.~21 (thermodynamique de l'atmosphère)~: enclume d'un orage supercellulaire, Stephen Corfidi, NOAA, domaine public~; radiosondages de l'Integrated Global Radiosonde Archive, I. Durre et al., NOAA NCEI, domaine public.
  \item Chap.~22 (systèmes météorologiques)~: une dépression au-dessus de l'Islande, NASA/GSFC, MODIS Rapid Response Team, Jacques Descloitres, domaine public~; l'œil de l'ouragan Isabel vu de la Station spatiale internationale, Ed Lu, NASA, domaine public.
  \item Chap.~23 (variabilité climatique)~: indices de l'oscillation nord-atlantique et ONI, NOAA Climate Prediction Center, domaine public.
  \item Chap.~24 (archives climatiques et indicateurs)~: couches annuelles de la carotte de glace GISP2 à 1855~m, Anthony Gow, U.S. Army Corps of Engineers CRREL (NOAA Paleoclimatology), domaine public~; coupe d'un spéléothème de travertin, Diamond Caverns, Kentucky, James St.~John, CC BY 2.0~; la compilation benthique LR04 de L.\,E. Lisiecki et M.\,E. Raymo (2005), l'enregistrement de température d'EPICA Dome C de J. Jouzel et al. (2007) et la compilation antarctique de \ce{CO2} de B. Bereiter et al. (2015), NOAA NCEI Paleoclimatology, domaine public.
  \item Chap.~25 (forçage orbital et âges glaciaires)~: éléments orbitaux et insolation de Berger et Loutre (1991), NOAA NCEI Paleoclimatology, domaine public.
  \item Chap.~26 (le climat à travers les temps géologiques)~: enregistrement isotopique benthique cénozoïque CENOGRID, T. Westerhold et al. (2020), PANGAEA, CC BY 4.0~; compilation phanérozoïque du \ce{CO2} de G.\,L. Foster, D.\,L. Royer et D.\,J. Lunt (2017), \emph{Nature Communications}, CC BY 4.0.
  \item Chap.~27 (changement climatique anthropique)~: sixième rapport d'évaluation du GIEC, groupe de travail~I, données étendues de l'annexe~III, CC BY 4.0.
  \item Chap.~28 (aléas géologiques et risque)~: dépôts de lave torrentielle à Caraballeda, Venezuela, Lawson Smith, U.S. Army Corps of Engineers, domaine public~; une mosquée restée debout après le tsunami de 2004 à Aceh, USGS, domaine public.
  \item Chap.~29 (gisements de minerai)~: la mine de Bingham Canyon vue de la Station spatiale internationale, NASA (équipage de l'Expedition~15), domaine public~; chapeau de fer de la mine Flambeau, Wisconsin, James St.~John, CC BY 2.0~; données teneur--tonnage de D.\,A. Singer, V.\,I. Berger et B.\,C. Moring (2008), USGS, domaine public.
  \item Chap.~30 (géologie de l'énergie)~: centrale géothermique près de Geyserville, Californie, U.S. Department of Defense, domaine public~; masses volumiques du dioxyde de carbone du NIST Chemistry WebBook (Standard Reference Database 69).
  \item Chap.~31 (méthodes de terrain et de laboratoire~II)~: augite poecilitique englobant du plagioclase, Strekeisen, CC BY-SA 4.0~; lame mince d'un micaschiste à grenat, Jackdann88, CC BY 3.0~; microsonde électronique du Service géologique de Finlande, Kallerna, CC BY-SA 4.0~; teneurs en silice de matériaux de référence rocheux, J.\,E. Taggart et D.\,F. Siems (2002), USGS, domaine public.
\end{itemize}
